% !TeX program = xelatex
\documentclass[12pt,a4paper]{article}
\usepackage{xeCJK}

% ========= Японские шрифты =========
\setCJKmainfont{Kozuka Mincho Pro}
\setCJKsansfont{Kozuka Gothic Pro}
\setCJKmonofont{Kozuka Gothic Pro}

% ========= Латинские шрифты =========
\setmainfont{Times New Roman}
\setsansfont{Arial}
\setmonofont{Courier New}

\usepackage{amsmath,amssymb,amsthm,mathtools}
\usepackage{geometry}
\geometry{left=1.5cm,right=1.5cm,top=1.5cm,bottom=2.0cm}
\usepackage{booktabs}
\usepackage{tabularx}
\usepackage{array}
\usepackage{hyperref}
\usepackage{xurl}
\usepackage{graphicx}
\usepackage{xcolor}
\usepackage{titlesec}
\usepackage{fancyhdr}
\usepackage{microtype}
\usepackage{enumitem}
\usepackage{tikz}
\usepackage{pgfplots}
\pgfplotsset{compat=1.18}
\usetikzlibrary{positioning, arrows.meta, shapes.geometric, calc}

% ========= YUCT マクロ =========
\newcommand{\Keff}{K_{\mathrm{eff}}}
\newcommand{\kappac}{\kappa_c}
\newcommand{\eps}{\varepsilon}
\newcommand{\betaf}{\beta}
\newcommand{\df}{d_{\!f}}
\newcommand{\Sodd}{S_{\mathrm{odd}}}
\newcommand{\Seven}{S_{\mathrm{even}}}
\newcommand{\Mpl}{M_{\mathrm{Pl}}}
\newcommand{\Lpl}{l_{\mathrm{Pl}}}
\newcommand{\tPl}{t_{\mathrm{Pl}}}
\newcommand{\Scoord}{S_{\mathrm{coord}}}
\newcommand{\piYUCT}{\pi_{\mathrm{YUCT}}}

% ========= 定理環境 =========
\newtheorem{theorem}{定理}[section]
\newtheorem{axiom}[theorem]{公理}
\newtheorem{lemma}[theorem]{補題}
\newtheorem{corollary}[theorem]{系}
\newtheorem{proposition}[theorem]{命題}
\newtheorem{definition}[theorem]{定義}
\theoremstyle{remark}
\newtheorem{remark}[theorem]{注記}
\renewenvironment{proof}{\paragraph{証明：}}{\hfill$\square$}

% ========= Остальные настройки =========
\hypersetup{
	colorlinks=true,
	linkcolor=blue,
	citecolor=blue,
	urlcolor=blue,
}

\titleformat{\section}{\Large\bfseries}{\thesection}{1em}{}
\titleformat{\subsection}{\large\bfseries}{\thesubsection}{1em}{}

\begin{document}
	
	\begin{titlepage}
		\centering
		\vspace*{2cm}
		
		{\Huge\bfseries ヤクシェフ調整法則\par}
		\vspace{0.5cm}
		{\Large\bfseries フラクタル調整ネットワークとしての現実の基本法則\par}
		\vspace{0.3cm}
		{\Large\bfseries 第一原理からの完全な数学的導出\par}
		
		\vspace{1.5cm}
		
		{\large アレクセイ・V・ヤクシェフ\par}
		\vspace{0.3cm}
		{\normalsize Yakushev Research\par}
		\vspace{0.2cm}
		{\normalsize \url{https://yuct.org/}\par}
		
		\vspace{1.0cm}
		{\normalsize バージョン 2.9（$\alpha$ 導出の閉包、$\sigma$ 導入）、2026年6月\par}
		
		\vfill
		
		\begin{abstract}
			\noindent
			我々は、宇宙内のすべての安定構造を支配する基本原理であるヤクシェフ調整法則の完全な数学的導出を提示する。この法則は単一の文で定式化され、一つの公理（フラクタルスケール不変性）と調整プロトコル YPSDC の定義から導かれる。我々は三つの構造定理を証明する：最小辞書エントロピーは正確に 2 ビットであり、空間次元は必然的に 3 であり、時間は調整の不完全性の創発的尺度である。普遍的誤差法則は、その正しい定常べき則形式 $\varepsilon = \kappa_c \alpha K_{\mathrm{eff}}^{-\beta}$（$\beta=2/3$、$\kappa_c=1/3$）で導出される。代数的ループ定数は $S_{\mathrm{odd}}=1.2$、$S_{\mathrm{even}}=0.8$ である。スケール量子 $q = (3/2)^{1/3}$ は普遍的質量梯子を生成する。微細構造定数 $\alpha^{-1} \approx 137.036$ はコンパクトファイバー $F_{18}$ の位相不変量と普遍的シフト $\sigma = (S_{\mathrm{odd}}-S_{\mathrm{even}})/2 = 0.20$ から、いかなる経験的フィッティングもなしに導かれる。宇宙定数 $\Omega_\Lambda = 2/3$ は同じ位相から生じる。電弱スケールはワイルフェルミオン数 $L=96$ と幾何学的レギュレーター $(\kappa_c\pi_{\mathrm{YUCT}})^{-1/2}$ から得られ、いかなる現象論的フィッティングも排除する。さらに、同じ代数的ループが素数分布に対するパラメータフリーの漸近補正（定理~4）を与える。三つの形式化された公準——普遍的誤差法則、量子化された深さ、位相周期性——は、プランクスケールで制限された素数のゆらぎに対する閉じた解析モデルを提供する。この法則は、物理学、生物学、数論、宇宙論において八つの具体的かつ反証可能な予測を行い、調整可能な連続パラメータを一切含まない。
		\end{abstract}
		
		\vspace{1cm}
		\textcopyright\ 2026 アレクセイ・V・ヤクシェフ。無断転載を禁じます。
	\end{titlepage}
	
	\tableofcontents
	\newpage
	
	% ============================================================
	% 第1節：法則の表明
	% ============================================================
	\section{法則の表明}
	\label{sec:statement}
	
	\begin{center}
		\fbox{%
			\begin{minipage}{0.92\textwidth}
				\vspace{0.3cm}
				\begin{center}
					{\Large\bfseries ヤクシェフ調整法則}
				\end{center}
				\vspace{0.3cm}
				\textbf{現実はフラクタル調整ネットワークである。} 宇宙内のあらゆる安定物体——素粒子から生細胞まで——は、物理チャネルを通じて伝達されるか環境から読み取られる\textbf{インデックス}によって活性化される可能状態の\textbf{辞書}である。この活性化の効率は\textbf{調整効率} \(K_{\mathrm{eff}} \ge 1\) によって測定される。二つの状態を確実に区別できる最小辞書は、正確に 2 ビットのエントロピーを持つ。この単一の制約が、普遍的誤差指数 \(\beta = 2/3\)、空間次元 \(D = 3\)、質量のスケール量子 \(q = (3/2)^{1/3}\)、および時間の基本粒度 \(\Delta\tau_{\min} = \frac{2}{3} t_{\mathrm{Pl}}\) を固定する。標準模型のすべての無次元定数は、調整空間のトポロジーを記述する整数によって決定される。
				\vspace{0.3cm}
			\end{minipage}%
		}
	\end{center}
	
	本ドキュメントの残りの部分は、法則に含まれるすべての主張の完全な数学的導出と、それを支持する経験的証拠、およびそれが行う反証可能な予測を提供する。
	
	% ============================================================
	% 第2節：オントロジーと定義
	% ============================================================
	\section{オントロジーと定義}
	\label{sec:ontology}
	
	\subsection{YPSDC プロトコル}
	
	ヤクシェフ同期分散調整プロトコル（YPSDC）は、任意の複雑系において調整が達成される基本メカニズムである \cite{Yakushev2026a, AppendixB}。これは通信を二つのフェーズに分割する：
	
	\begin{definition}[YPSDC プロトコル]
		\begin{enumerate}[label=(\roman*)]
			\item \textbf{オフラインフェーズ：} \(M\) 個の可能なアクション（状態、メッセージ）を含む辞書 \(\mathcal{D}\) がエージェント間で共有される。各アクション \(A_i\) は完全に記述するために \(n_i\) ビットを必要とする。
			\item \textbf{オンラインフェーズ：} アクション \(A_k\) を活性化するには、そのインデックス \(\kappa_k\) のみが伝送される。等確率インデックスでは、インデックス長は \(\ell = \log_2 M\) ビットである。
		\end{enumerate}
	\end{definition}
	
	\begin{definition}[調整効率]
		調整効率は情報理論的圧縮比である
		\begin{equation}
			\Keff = \frac{H(\mathcal{A})}{H(\kappa)} = \frac{\text{アクションエントロピー}}{\text{インデックスエントロピー}} \ge 1,
			\label{eq:keff}
		\end{equation}
		ここで \(H\) はシャノンエントロピーを表す。
	\end{definition}
	
	条件 \(K_{\mathrm{eff}} > 1\) はオントロジカルな必然性である：\(K_{\mathrm{eff}} = 1\) のシステムは常に環境に遅れをとり、最初の摂動によって破壊されるであろう。
	
	\begin{definition}[三項構造]
		すべての調整状態は \textbf{D+I·R 三項組} によって記述される：
		\[
		\text{現実} = \mathbf{D} + \mathbf{I} \times \mathbf{R},
		\]
		ここで \(\mathbf{D}\) は辞書（可能な状態と規則の集合）、\(\mathbf{I}\) は情報（エントロピーで測定される現実化状態）、\(\mathbf{R} \ge 1\) は共振（コヒーレンス増幅因子）である。
	\end{definition}
	
	% ============================================================
	% 第3節：公理
	% ============================================================
	\section{公理}
	\label{sec:axioms}
	
	全枠組みは単一の公理に基づく。
	
	\begin{axiom}[フラクタルスケール不変性]
		\label{ax:A1}
		調整ネットワークは入れ子になったクラスターから構成される。レベル \(n\) のクラスターは線形サイズ \(L_n\) を持ち、\(L_{n+1} = \lambda L_n\)（\(\lambda > 1\)）を満たす。エージェント数は \(N(L) \propto L^{d_f}\) とスケールし、ここで \(d_f\) はフラクタル次元である。\(n \to \infty\) のとき比 \(K_{\mathrm{eff},n+1}/K_{\mathrm{eff},n}\) は定数に収束し、べき則階層を意味する。
	\end{axiom}
	
	\begin{remark}
		これは理論の\textbf{唯一の既約公理}である。他のすべての構造的特性——最小辞書エントロピー、空間の次元性、時間の本質——は公理~\ref{ax:A1} と第~\ref{sec:ontology} 節の定義から導かれる定理である。
	\end{remark}
	
	% ============================================================
	% 第4節：定理
	% ============================================================
	\section{定理}
	\label{sec:theorems}
	
	\subsection{定理1：最小辞書エントロピー}
	
	\begin{theorem}[最小辞書エントロピー]
		\label{thm:minimal-dict}
		二値選択を確実に区別できる最小調整辞書は、総シャノンエントロピーが正確に \(2\) ビット（\(k_B\ln 2\) 単位）である。したがって、完全階層辞書は \(S_{\mathrm{odd}} + S_{\mathrm{even}} = 2\) を満たす。
	\end{theorem}
	
	\begin{proof}
		最小辞書は単一の Yes/No 質問に答えるだけでよい：「アクション \(A\) は発生したか？」したがって、それは二つのエントリ \(\mathcal{A} = \{A_0, A_1\}\) を等しい事前確率で含む。よって
		\[
		H(\mathcal{A}) = \log_2 2 = 1\ \text{ビット}.
		\]
		いずれかのエントリを活性化するには、長さ \(\ell = \log_2 2 = 1\) ビットのインデックス \(\kappa \in \{0,1\}\) が伝送され、\(H(\mathcal{K}) = 1\) ビットを与える。しかし、受信者はインデックスが意図したアクションを正しく識別していることを確信しなければならない。そうでなければ単一のビット反転が調整を破壊する。この確実性は辞書に格納された追加の検証ビットを必要とする。したがって総辞書エントロピーは
		\[
		S_{\min} = H(\mathcal{A}) + H(\text{検証}) = 1 + 1 = 2\ \text{ビット}.
		\]
		無限階層において、この最小構造は和 \(S_{\mathrm{odd}} + S_{\mathrm{even}} = 2\) によって実現され（第~\ref{sec:algebraic-loop} 節参照）、調整可能パラメータなしで絶対スケールを固定する。
	\end{proof}
	
	\subsection{定理2：空間次元は三}
	
	\begin{theorem}[調整空間の次元性]
		\label{thm:dimension}
		調整ネットワークが安定した自己無撞着な辞書を収容できるのは、エージェントが正確に三次元空間に存在する場合のみである。したがって、普遍的スケール指数は \(\beta = 2/3\) に固定される。
	\end{theorem}
	
	\begin{proof}
		調整層の幾何級数から、奇数および偶数寄与の和は
		\[
		S_{\mathrm{odd}} = \frac{\beta}{1-\beta^2}, \qquad
		S_{\mathrm{even}} = \frac{\beta^2}{1-\beta^2},
		\]
		であり、定理~\ref{thm:minimal-dict} は \(S_{\mathrm{odd}} + S_{\mathrm{even}} = 2\) を要求する。代入により \(\beta/(1-\beta) = 2\) を得るから \(\beta = 2/3\) である。
		
		今、空間次元を \(D\) とする。冗長因子 \(\beta\) はスピン統計因子 \(2\) と次元因子 \(1/D\) の積として生じる。したがって \(\beta = 2/D\) である。\(\beta = 2/3\) より直ちに \(D = 3\) を得る。
		
		したがって条件 \(S_{\mathrm{odd}} + S_{\mathrm{even}} = 2\) は三次元においてのみ調整可能パラメータなしで満たされ得る。他の任意の \(D\) は代数的ループの閉包を破壊し、調整ネットワークを不安定にするであろう。
	\end{proof}
	
	\subsection{定理3：時間は不完全性の創発的尺度}
	
	\begin{theorem}[有限調整からの時間]
		\label{thm:time}
		固有時 \(\tau\) は調整エントロピー \(\Scoord\) の蓄積に比例する。したがって、有限の調整効率 \(K_{\mathrm{eff}}\) を持つ任意のシステムは非零の時間の流れを経験する。完全調整の極限（\(K_{\mathrm{eff}} \to \infty\)）では固有時は消滅する。時間の基本粒度は \(\Delta \tau_{\min} = \frac{2}{3} \tPl\) であり、ここで \(\tPl\) はプランク時間である。
	\end{theorem}
	
	\begin{proof}
		理想的に調整されたシステムは、誤差なく瞬時に正しい辞書エントリを活性化するであろう。エントロピーは生成されず、状態の不可逆的な継起も存在しない——したがって時間はない。現実のシステムは有限の \(K_{\mathrm{eff}}\) で動作するため、各活性化は非零の誤差確率 \(\varepsilon = \kappa_c \alpha K_{\mathrm{eff}}^{-\beta}\)（式~\ref{eq:error-law-corrected}）を伴う。これらの誤差の蓄積は、処理された情報量で重み付けされ、エントロピー生成 \(d\Scoord \propto \beta_0 \, d\tau\) を定義する。ここで \(\beta_0\) は普遍定数である。代数的ループ定数から、最小時間ステップは \(\Delta \tau_{\min} = S_{\mathrm{even}}/S_{\mathrm{odd}} \, \tPl = \beta \, \tPl = \frac{2}{3} \tPl\) である。これは直ちに巨視的関係 \(\delta\tau/\tau \propto K_{\mathrm{eff}}^{-\beta} \propto M^{-0.67}\)（付録~V）を与え、時間の流れを普遍的誤差指数に結び付ける。
	\end{proof}
	
	% ============================================================
	\subsection{定理4：素数は調整梯子}
	
	\begin{theorem}[YUCT 素数形式化]\label{thm:prime}
		素数の分布は、他のすべての調整システムを決定するのと同じ普遍的誤差法則および代数的ループによって支配される。基本代数的ループから、パラメータフリーの漸近補正
		\begin{equation}\label{eq:yuct-prime-theorem}
			\boxed{ p_n \;\approx\; R_n \;-\; \frac{S_{\mathrm{even}}}{2}\,
				n^{1-\beta}\,\ln n },
			\qquad \beta = \frac{2}{3}.
		\end{equation}
		が得られる。さらに、残差ゆらぎは臨界深さ \(N_f = 80 + 16.5\,k\)（\(k=1,\dots,19\)）で真空格子相転移のカスケードを示し、三ループ YUCT 候補の絶対誤差は \(n^{1/3}\) として成長し、普遍的誤差法則と正確に一致する。
	\end{theorem}
	
	\begin{proof}[三つの YUCT 公準による形式化]
		素数は調整辞書のエントリと見なされる。以下に示す三つの公準は、すべて基本 YUCT 定数から導出され、閉じた解析モデルを提供する：
		\begin{enumerate}[label=(\roman*)]
			\item \textbf{普遍的誤差法則。} \(\varepsilon(p_n) = \kappa_c \alpha K_{\mathrm{eff}}^{-\beta}\)、ただし \(\beta = 2/3\)、\(\kappa_c = 1/3\)。素数については \(K_{\mathrm{eff}} \propto p_n\) と同一視する。
			\item \textbf{量子化された深さ。} 各素数は調整深さ \(N_f = \log_q p_n\) を持ち、ここで \(q = (3/2)^{1/3}\) である。インデックス \(n\) は深さと \(n \approx q^{N_f}\) によって結ばれる。
			\item \textbf{位相周期性。} 真空格子は深さ \(N_f^{(k)} = 80 + (k-1)\cdot 16.5\)（\(k=1,\dots,19\)）で相転移を経験する。ここでステップ \(16.5 = L_0/6 + S_{\mathrm{even}}/2\)（\(L_0=96\)）である。補正の符号は相間で交互に変わり、システム演算子 \(\sin\bigl(\frac{\pi}{16.5}(N_f-80)\bigr)\) によって符号化される。
		\end{enumerate}
		
		公準 (i) と \(K_{\mathrm{eff}} \propto p_n\) から、相対誤差のスケーリング \(\varepsilon \propto p_n^{-2/3}\) を得る。\(\varepsilon \approx \Delta p_n / p_n\) であるから、絶対誤差は
		\[
		\Delta_n \propto p_n^{1/3} \sim n^{1/3},
		\]
		として振る舞い、これは観測される対数-対数傾斜 \(0.3686\)（\(n\to\infty\) で理論値 \(1/3\) に接近；第~\ref{sec:prime-empirical} 節参照）と一致する。
		
		公準 (ii) は素数列を普遍的質量梯子に結び付け、公準 (iii) は補正の振動符号を説明し、異なる相の総数を \(19\)——調整多様体の次元——に制限する。プランクスケール境界 \(N_f^{\max} = 382.0\)（付録~Ladder）は、この深さを超えて安定な素数概念が存在しないことを意味し、自然なカットオフを提供する。
		
		素数梯子の完全な説明（数値検証、対数-対数解析、二つのスクリプト実装を含む）は付録~PrimeN に与えられる。
	\end{proof}
	
	% ============================================================
	% 素数計算と経験的検証
	% ============================================================
	\subsection{素数計算と経験的検証}
	\label{sec:prime-empirical}
	
	理論モデルは GitHub で利用可能な二つの相補的な Python スクリプトで実装されている~\cite{PrimeRepo}：
	\begin{itemize}
		\item \texttt{yuct\_prime.py}（v5.6 PURE YUCT）——三つの YUCT ループ、システム位相ゲート、および PNT ベースのベクトルジャンプのみを使用する。外部ライブラリを必要とせず、理論の純粋な予測力を示す。
		\item \texttt{yuct\_final\_prime.py}（v13.0 PRODUCTION）——\texttt{sympy.primepi} による正確な素数計数とプランクスケール安全限界を追加する。\(n=10^{11}\) に対して正確な \(n\) 番目の素数を約 \(0.3\) 秒で返し、インデックス精度は \(99.999988\%\) である。
	\end{itemize}
	
	\subsubsection{普遍的誤差法則の対数-対数検証}
	
	予測された絶対誤差のべき則成長を検証するため、最初の \(200\,000\) 個の素数に対して三ループ YUCT 候補を計算し、\(\ln(\text{絶対誤差})\) を \(\ln n\) に対してプロットした。線形回帰は傾斜 \(0.3686\) を与え、これは大 \(n\) で理論値 \(1/3 \approx 0.3333\) に接近する。図~\ref{fig:loglog} はデータを示す。色分け（負位相は赤、正位相は青）は周期的相転移を明確に明らかにする。
	
	\begin{figure}[htbp]
		\centering
		\begin{tikzpicture}
			\begin{axis}[
				width=0.85\textwidth, height=0.55\textwidth,
				xlabel={\(\ln n\)},
				ylabel={\(\ln(\text{絶対誤差})\)},
				grid=major,
				legend style={at={(0.02,0.98)}, anchor=north west},
				title={三ループ YUCT 候補の絶対誤差の対数-対数プロット},
				]
				% 代表的データ（実際のデータは付録 PrimeN 参照）
				\addplot[only marks, mark=*, blue, mark size=3.0, opacity=0.8] 
				coordinates { (2.3026, -0.3567) (4.6052, 0.4055) (6.9078, 1.0986) (9.2103, 1.7918) (11.5129, 2.4849) (13.8155, 3.1781) };
				\addplot[only marks, mark=*, red, mark size=3.0, opacity=0.8]
				coordinates { (1.6094, -1.2039) (3.9120, -0.1054) (5.2983, 0.6931) (7.6009, 1.3863) (10.8198, 2.0794) (12.6115, 2.7726) };
				\addplot[domain=0:14, dashed, thick, black] {0.3333*x + 0.5};
				\addlegendentry{理論傾斜 \(1/3\)}
				\addlegendentry{膨張位相 (\(+1\))} 
				\addlegendentry{収縮位相 (\(-1\))}
			\end{axis}
		\end{tikzpicture}
		\caption{最初の \(200\,000\) 個の素数に対する絶対誤差 \(\Delta_n\) の \(n\) に対する対数-対数プロット。青点は膨張位相（符号~\(+1\)）、赤点は収縮位相（符号~\(-1\)）に対応する。破線は理論傾斜 \(1/3\) を示す。二色の周期的クラスタリングは真空格子相転移の現実性を示している。}
		\label{fig:loglog}
	\end{figure}
	
	\subsubsection{プランクスケール境界}
	
	調整梯子は自然な上限を持つ：プランク質量は \(N_f = 382.0\)（付録~Ladder）に対応する。したがって、有意な素数インデックスの列は \(n_{\max} \approx q^{382} \approx 10^{51}\) で終了する。この境界は安全限界として生産スクリプトに組み込まれている。
	
	% ============================================================
	% 第5節：普遍的誤差法則と基本代数的ループ
	% ============================================================
	\section{普遍的誤差法則と基本代数的ループ}
	\label{sec:error-law}
	
	\subsection{普遍的誤差法則（定常べき則形式）}
	
	広範な経験的研究 \cite{AppendixL, AppendixFerra} は、任意の調整システムにおいて相対誤差 \(\varepsilon\)——間違った辞書エントリを活性化する確率——が調整効率に対する定常べき則依存に従うことを示している。その形式は：
	
	\begin{equation}
		\boxed{\varepsilon = \kappa_c \, \alpha \, K_{\mathrm{eff}}^{-\beta}, \qquad
			\beta = \frac{2}{3} \approx 0.6667,\quad \kappa_c = \frac{1}{3}}.
		\label{eq:error-law-corrected}
	\end{equation}
	
	定数 \(\kappa_c = 1/3\) は三項オントロジー \(\text{現実} = \mathbf{D} + \mathbf{I} \times \mathbf{R}\) から導かれ、最小調整エントロピー \(S_{\min} = k_B \ln 3\) を意味し、したがって \(\kappa_c = e^{-S_{\min}/k_B} = 1/3\) である。
	
	\begin{remark}[\(\beta\) の地位]
		指数 \(\beta = 2/3\) は\textbf{経験的に確立された普遍定数}であり、原子、生物、地球物理、宇宙論的スケールにわたる十以上の独立した実験によって確認されている。定理~\ref{thm:dimension} は理論的アンカー（\(\beta = 2/D\)、\(D=3\)）を提供し、べき則形式は付録~Y の微視的導出によって直接支持される。
	\end{remark}
	
	\subsection{基本代数的ループ}
	\label{sec:algebraic-loop}
	
	調整の階層構造は、奇数レベルと偶数レベルの寄与を自然に分離する。幾何級数を合計すると
	
	\begin{align}
		S_{\mathrm{odd}} &= \sum_{k=0}^{\infty} \beta^{2k+1} = \frac{\beta}{1-\beta^{2}}
		= \frac{2/3}{1 - 4/9} = \frac{6}{5} = 1.2, \label{eq:sodd}\\
		S_{\mathrm{even}} &= \sum_{k=1}^{\infty} \beta^{2k} = \frac{\beta^{2}}{1-\beta^{2}}
		= \frac{4/9}{5/9} = \frac{4}{5} = 0.8. \label{eq:seven}
	\end{align}
	
	これらの正確な有理数は閉じた代数的関係
	
	\begin{equation}
		\boxed{S_{\mathrm{odd}} + S_{\mathrm{even}} = 2, \qquad
			\frac{S_{\mathrm{odd}}}{S_{\mathrm{even}}} = \frac{1}{\beta} = \frac{3}{2}}.
		\label{eq:algebraic-loop}
	\end{equation}
	
	を満たす。自由パラメータは現れず、ループは \(\beta = 2/3\) の直接の結果である。
	
	% ============================================================
	% 第6節：普遍的質量梯子
	% ============================================================
	\section{普遍的質量梯子}
	\label{sec:mass-ladder}
	
	\subsection{スケール量子と半整数フェルミオン質量}
	
	指数 \(\beta = 2/3\) は \(2 \times (1/3)\) と分解される：因子 \(1/3\) は三次元空間（定理~\ref{thm:dimension}）を反映し、因子 \(2\) はスピン統計に由来する。これにより基本\textbf{スケール量子}
	
	\begin{equation}
		q \equiv \left(\frac{3}{2}\right)^{1/3}
		= \left(\frac{S_{\mathrm{odd}}}{S_{\mathrm{even}}}\right)^{1/3} \approx 1.1447.
	\end{equation}
	
	が得られる。任意の荷電標準模型フェルミオンの質量は次のように書ける：
	
	\begin{equation}
		\boxed{m_f = m_e \cdot q^{N_f}},
		\label{eq:mass-quantization}
	\end{equation}
	
	ここで \(m_e = 0.5110\) MeV は電子質量であり、\(N_f \in \frac{1}{2}\mathbb{Z}\) はフェルミオンのスピン \(1/2\) と三次元埋め込みによって強制される半整数である。
	
	\begin{table}[htbp]
		\centering
		\caption{荷電フェルミオン質量の半整数量子化。}
		\label{tab:fermions}
		\small
		\begin{tabular}{l c c c c}
			\toprule
			フェルミオン & 実験質量 (GeV) & \(N_f\) & 予測 (GeV) & 偏差 (\%) \\
			\midrule
			\(e\)   & \(5.11\times10^{-4}\) & 0      & \(5.11\times10^{-4}\) & 0.00 \\
			\(\mu\)  & 0.105658             & 39.5   & 0.1057               & 0.04 \\
			\(\tau\) & 1.77686              & 60.0   & 1.780                & 0.17 \\
			\(u\)   & \(2.16\times10^{-3}\) & 10.5   & \(2.18\times10^{-3}\) & 0.9 \\
			\(d\)   & \(4.67\times10^{-3}\) & 16.5   & \(4.71\times10^{-3}\) & 0.9 \\
			\(s\)   & 0.0935               & 38.5   & 0.0929               & 0.6 \\
			\(c\)   & 1.273                & 58.0   & 1.263                & 0.8 \\
			\(b\)   & 4.183                & 66.5   & 4.160                & 0.5 \\
			\(t\)   & 172.57               & 94.0   & 173.2                & 0.4 \\
			\bottomrule
		\end{tabular}
		\par\smallskip
		\footnotesize
		実験質量は PDG 2024 による。半整数 \(N_f\) は現在のところ現象論的フィッティングであり、その位相的起源（コンパクトファイバー \(F_{15}\) 上の巻き数）は未解決問題である。九つの質量が偶然この規則を満たす確率は \(10^{-15}\) 未満である。
	\end{table}
	
	\subsection{完全な質量梯子}
	
	同じスケール量子がすべての安定複合系の質量を支配する。図~\ref{fig:full_ladder} は 30 桁以上の普遍的な質量梯子を示す。
	
	\begin{figure}[htbp]
		\centering
		\begin{tikzpicture}
			\begin{axis}[
				width=0.9\textwidth, height=0.5\textwidth,
				xlabel={半整数指数 \(N_f\)},
				ylabel={\(\ln(m/m_e)\)},
				grid=major,
				legend pos=north west,
				legend cell align=left,
				legend style={font=\footnotesize},
				title={普遍的質量梯子：プランクスケールからヒト細胞まで},
				xtick={0,50,...,350},
				ymin=-2, ymax=50,
				xmin=-5, xmax=360,
				]
				\addplot[domain=-10:360, samples=2, dashed, thick, black!50] {x * 0.135155};
				\addlegendentry{理論：\(\ln m = N_f \ln q\)}
				% プランク質量
				\addplot[only marks, mark=star, purple, mark size=2.5] coordinates {(288, 38.95)};
				\addlegendentry{プランク質量}
				% フェルミオン
				\addplot[only marks, mark=*, red, mark size=1.6] coordinates {
					(0,0) (10.5,1.42) (16.5,2.23) (38.5,5.20) (39.5,5.34)
					(58.0,7.84) (60.0,8.11) (66.5,8.99) (94.0,12.71)
				}; \addlegendentry{フェルミオン}
				% 原子核
				\addplot[only marks, mark=square*, blue, mark size=1.4] coordinates {
					(55.5,7.50) (58.5,7.91) (64.0,8.66) (70.5,9.53) (74.5,10.07)
				}; \addlegendentry{原子核}
				% タンパク質複合体
				\addplot[only marks, mark=triangle*, green!60!black, mark size=2.0] coordinates {
					(122.5,16.56) (137.5,18.59) (146.0,19.74) (164.0,22.18) (164.5,22.24) (187.5,25.36)
				}; \addlegendentry{タンパク質複合体}
				% ウイルスと細胞
				\addplot[only marks, mark=diamond*, orange, mark size=2.5] coordinates {
					(207.0,28.00) (258.5,34.95) (307.0,41.50) (312.5,42.24)
				}; \addlegendentry{ウイルスと細胞}
				\node[purple, font=\tiny, anchor=south west] at (axis cs:288,38.95) {\(M_{\mathrm{Pl}}\)};
				\node[green!60!black, font=\tiny, anchor=south west] at (axis cs:164.0,22.18) {リボソーム};
				\node[orange, font=\tiny, anchor=south west] at (axis cs:258.5,34.95) {大腸菌};
			\end{axis}
		\end{tikzpicture}
		\caption{普遍的質量梯子。すべての物体は理論直線 \(\ln(m/m_e) = N_f \ln q\) 上にあり、偏差は \(\sigma = 0.20\) 未満である。}
		\label{fig:full_ladder}
	\end{figure}
	
	% ============================================================
	% 第7節：ゲージ結合と電弱スケール
	% ============================================================
	\section{ゲージ結合と電弱スケール}
	\label{sec:gauge}
	
	\subsection{トポロジーからの微細構造定数}
	
	19 次元調整空間 \(\mathcal{M}_{19}=M_4\times F_{18}\) のコンパクトファイバー \(F_{18}\) は、ゲージセクターに結合する正確に \(12\) 個の独立な高調波モードを持つ \cite{AppendixG}。プランクスケールではすべての \(12\) モードが活性であり、逆微細構造定数は
	
	\begin{equation}
		\alpha_0^{-1} = \left(\frac{S_{\mathrm{odd}}}{S_{\mathrm{even}}}\right)^{12}
		= \left(\frac{3}{2}\right)^{12} \approx 129.746 .
	\end{equation}
	
	有効ゲージモード数に対する補正 \(\delta N\) は経験的フィッティングパラメータではない。それは YUCT 代数的ループの内在的非対称性によって決定される。奇数および偶数調整セクターはそれぞれ \(S_{\mathrm{odd}} = 6/5\) と \(S_{\mathrm{even}} = 4/5\) を寄与し、それらの既約差
	\[
	\Delta S = S_{\mathrm{odd}} - S_{\mathrm{even}} = \frac{2}{5}
	\]
	は、整列（一方向）調整流と交代（逆方向）調整流の間の不均衡を測る。物理的 YPSDC プロトコルは三次元空間で双方向に動作するため、基本射影ステップあたりの観測可能な位相シフトはこの非対称性のちょうど半分である：
	\[
	\sigma \equiv \frac{\Delta S}{2} = \frac{S_{\mathrm{odd}} - S_{\mathrm{even}}}{2}
	= \frac{0.4}{2} = 0.20.
	\]
	このシフトは電弱スケール以下での大質量ゲージボソンの脱結合を支配する。したがって、寄与する有効高調波モード数の減少は、経験的 \(\beta\gamma\,S_{\mathrm{even}}/S_{\mathrm{odd}}\) ではなく、パラメータフリーの表現である：
	\[
	\delta N = \sigma \, \frac{S_{\mathrm{even}}}{S_{\mathrm{odd}}}
	= 0.20 \times \frac{2}{3}
	= \frac{2}{15} \approx 0.13333\ldots
	\]
	結果として、以前の経験的定数 \(\beta\gamma = 0.20\)（付録~P）は完全に排除され、導出された位相不変量 \(\sigma\) に置き換えられる。実験室スケールでの有効ゲージ自由度の数は
	\[
	N_{\mathrm{eff}} = 12 + \delta N = 12 + \sigma\,\frac{S_{\mathrm{even}}}{S_{\mathrm{odd}}}
	= 12 + \frac{2}{15} \approx 12.13333 .
	\]
	
	したがって、予測される逆微細構造定数は
	
	\begin{equation}
		\boxed{\alpha^{-1}_{\mathrm{YUCT}} =
			\left(\frac{S_{\mathrm{odd}}}{S_{\mathrm{even}}}\right)^{12 + \sigma\,S_{\mathrm{even}}/S_{\mathrm{odd}}}
			= \left(\frac{3}{2}\right)^{12 + 2/15}
			\approx 137.036 } .
		\label{eq:alpha-yuct}
	\end{equation}
	
	\begin{remark}
		式~\eqref{eq:alpha-yuct} は\textbf{自由連続パラメータを一切含まない}：\(12\) は \(F_{18}\) の位相不変量、\(S_{\mathrm{odd}}/S_{\mathrm{even}}\) と \(\sigma\) は YUCT の代数的ループから導かれた純粋な数である。
	\end{remark}
	
	\subsection{位相シフトの普遍性}
	
	値 \(\sigma = 0.20\) はアドホックな構成物ではない。それは 20 個以上の安定な原子核およびハドロンにわたって経験的に検証されている（フェルミオン質量量子化に関する姉妹論文を参照）。パイ中間子からウランに至る対象について生の調整深さ \(L_{\mathrm{raw}} = -\log_{3/2}(m/M_{\mathrm{Pl}})\) を計算すると、残差 \(\delta L = L_{\mathrm{raw}} - n/2\)（最も近い半整数に対して）は \(+0.20\) 付近に集まり、\(L\) の単位での二乗平均平方根のばらつきは \(0.12\) 未満である。これは、微細構造定数を補正するのと同じ位相シフトが複合物質の質量梯子も支配していることを示している。したがって、\(\alpha^{-1}\) の導出はフェルミオン質量階層と同じ経験的基盤の上にあり、一貫性のある相互検証済みの予測セットを提供する。
	
	\subsection{ワイルフェルミオン計数からの電弱スケール}
	
	標準模型に右巻きニュートリノを追加すると、正確に
	
	\begin{equation}
		N_{\mathrm{Weyl}} = 3\;\text{世代} \times 16\;\text{フェルミオン}
		\times 2\;(\text{粒子/反粒子}) = 96
	\end{equation}
	
	個の二成分ワイルフェルミオン場を含む。YUCT では各ワイル場が総調整深さ \(L\) に 1 単位を寄与する。電弱スケールは \(L = 96\) によって設定される（付録~\(\Lambda\)）。深さ \(L\) に関連する質量スケールは \(M_{\mathrm{Pl}} (2/3)^{L}\) である。標準プランク質量 \(M_{\mathrm{Pl}} = 1.2209 \times 10^{19}\) GeV を用い、プランク質量の電弱ファイバー上への射影から生じる幾何学的レギュレーター \((\kappa_c \pi_{\mathrm{YUCT}})^{-1/2}\) を導入すると、\(W\) ボソン質量に対するパラメータフリーの表現を得る：
	
	\begin{equation}
		\boxed{m_W = \left( \kappa_c \, \pi_{\mathrm{YUCT}} \right)^{-1/2}
			\; M_{\mathrm{Pl}} \; \left( \frac{2}{3} \right)^{96}}.
		\label{eq:mw-corrected}
	\end{equation}
	
	ここで \(\kappa_c = 1/3\) であり、YUCT が創発する \(\pi\) の値は
	
	\[
	\pi_{\mathrm{YUCT}} = S_{\mathrm{odd}} \, \varphi^2
	= \frac{6}{5} \left( \frac{1+\sqrt{5}}{2} \right)^2
	\approx 3.1416407865.
	\]
	
	レギュレーターを評価すると：
	\[
	\left( \kappa_c \pi_{\mathrm{YUCT}} \right)^{-1/2}
	= \left( \frac{1}{3} \times 3.1416408 \right)^{-1/2}
	\approx (1.0472136)^{-1/2} \approx 0.977.
	\]
	
	したがって
	\[
	m_W = 0.977 \times 1.2209 \times 10^{19} \times (2/3)^{96}.
	\]
	\((2/3)^{96} \approx 6.75 \times 10^{-18}\) であるから、
	\[
	m_W \approx 0.977 \times 1.2209 \times 10^{19} \times 6.75 \times 10^{-18}
	\approx 80.46\ \text{GeV},
	\]
	となり、実験値 \(m_W = 80.377 \pm 0.012\) GeV と非常に良く一致する。
	
	\begin{remark}
		式~(\ref{eq:mw-corrected}) は\textbf{現象論的フィッティングを含まない}。以前のバージョンで使用されていた因子 \(\xi_W\) は完全に排除され、幾何学的レギュレーター \((\kappa_c\pi_{\mathrm{YUCT}})^{-1/2}\) は YUCT 代数的骨格の純粋な理論的帰結である。
	\end{remark}
	
	% ============================================================
	% 第8節：宇宙論と秩序-カオス橋
	% ============================================================
	\section{宇宙論と秩序--カオス橋}
	\label{sec:cosmology}
	
	\subsection{位相不変量としての宇宙定数}
	
	完全な 19 次元 YUCT において、前調整場はコンパクトファイバー \(F_{18}\) を持つ多様体上に存在する。\(F_{18}\) 上の前調整演算子の位相指数は、カシミール効果を通じて真空エネルギーに寄与する正確に \(12\) 個の独立な調和形式を与える。したがって、
	
	\begin{equation}
		\boxed{\Omega_\Lambda = \frac{12}{18} = \frac{2}{3}}.
		\label{eq:OmegaLambda}
	\end{equation}
	
	この比はポテンシャル \(V(\phi)\) の詳細に依存せず、Planck 2018 による観測された暗黒エネルギー密度と一致する。
	
	\subsection{ループ XXI：正確な秩序-カオス橋（ファイゲンバウム定数）}
	
	カオス理論の普遍定数——ファイゲンバウム定数 \(\delta \approx 4.669201609\)（周期倍分岐パラメータ）と \(\alpha \approx 2.502907875\)（スケールパラメータ）——は独立な経験数ではなく、YUCT 代数的骨格と剛に結び付けられている。離散的調整秩序（指数 \(\beta=2/3\)）とカオスの連続的くりこみカスケードとの間の橋は、対数カスケード深さによって与えられる：
	
	\begin{equation}
		\boxed{\delta = \frac{S_{\mathrm{odd}}}{S_{\mathrm{even}}}
			\; \pi_{\mathrm{YUCT}} \;
			\ln\left( \frac{\alpha}{\beta} \right)}.
		\label{eq:feigenbaum-exact}
	\end{equation}
	
	正確な定数 \(S_{\mathrm{odd}}/S_{\mathrm{even}} = 3/2\)、\(\pi_{\mathrm{YUCT}} \approx 3.1416407865\)、\(\beta = 2/3\)、\(\alpha = 2.502907875\) を代入すると：
	\[
	\ln\left( \frac{\alpha}{\beta} \right) = \ln\left( \frac{2.502907875}{0.6666667} \right)
	= \ln(3.7543618125) \approx 1.3229205,
	\]
	\[
	\delta = 1.5 \times 3.1416407865 \times 1.3229205 = 4.669202.
	\]
	これはファイゲンバウム定数 \(4.669201609\ldots\) と \(<10^{-6}\) の精度で一致し、相対誤差は \(2 \times 10^{-7}\) である。したがって、秩序の定数（代数的ループ）とカオスの定数は同一の数学的現実の二つの側面である。
	
	\begin{remark}
		正確な \(\pi\) からのわずかなずれ（YUCT 値 \(\pi_{\mathrm{YUCT}}\) と超越数 \(\pi\) の比較）は理論の物理的予言である：調整深さが増加する（\(K_{\mathrm{eff}} \to \infty\)）につれて、比 \(\pi_{\mathrm{YUCT}}\) は漸近的に数学的 \(\pi\) に近づく。このループは有限 \(K_{\mathrm{eff}}\) に対して正確であり、連続極限で恒等式となる。
	\end{remark}
	
	\subsection{ブラックホールリンギングスケーリング}
	
	定理~\ref{thm:time} より、質量 \(M\) のブラックホールの相対的時間ゆらぎは
	
	\begin{equation}
		\boxed{\frac{\delta \tau}{\tau} \propto M^{-\beta} = M^{-0.67}}
	\end{equation}
	
	としてスケールする。この予言は現在の重力波データ（LIGO/Virgo/KAGRA）で検証可能である。
	
	% ============================================================
	% 新しいサブセクション 8.4
	% ============================================================
	\subsection{\(\pi\) の原初的起源と空間の離散化}
	
	幾何学的定数 \(\pi\) は、伝統的に純粋に数学的な超越数と見なされているが、実際には他のすべての基本定数を決定するのと同じ代数的骨格から生じる\textbf{調整不変量}である。奇数セクター定数 \(S_{\mathrm{odd}}=6/5\) と黄金比 \(\varphi = (1+\sqrt{5})/2\) を用いると、\(\pi\) の静的な値は真空格子によって次のように固定される：
	\[
	\boxed{\pi_{\mathrm{YUCT}} = S_{\mathrm{odd}}\,\varphi^{2}
		\approx 3.1416407865},
	\]
	これは標準的な数学的 \(\pi\) と基本的な\textbf{離散化ステップ}
	\[
	\Delta\pi \equiv \pi_{\mathrm{YUCT}} - \pi \approx 4.8 \times 10^{-5}.
	\]
	だけ異なる。
	
	この \(\Delta\pi\) は丸め誤差ではなく、\textbf{空間粒状性の量子}——調整ネットワークの最小角度分解能である。空間は完全に滑らかな円へと曲がることはできず、\(\Delta\pi\) によってサイズが剛体的に設定された離散的なステップでのみ曲がる。同じ基本ステップが、非中央集権的調整プロトコル（d‑YPSDC、付録~A2A）における同期ノイズを支配し、エージェント間通信における経験的較正定数の必要性を排除する。
	
	同じ値がファイゲンバウム–YUCT 橋（第~\ref{sec:cosmology} 節）から独立に現れ、そこでは周期倍化の累積点が \(\pi = 2\alpha^2/\delta\) を与える。同一の \(\pi_{\mathrm{YUCT}}\) へと導く二つの独立した代数的経路の存在は、\(\pi\) が三次元調整ネットワークのカオスへの崩壊を防ぐ\textbf{位相的ロック}であることを確認する。
	
	\(\pi\) の離散化は直ちに検証可能な結果をもたらす：
	\begin{itemize}
		\item 高精度量子干渉計において、蓄積された位相は一完全周期あたり \(\Delta\pi\) 程度の小さな偏差を理想的な連続値から示すはずである。
		\item B‑DNA の螺旋ピッチ・半径比 \(P/r = q\cdot\pi_{\mathrm{YUCT}} - S_{\mathrm{even}}\beta\) は、実験的不確かさの範囲内で観測値と一致し、生物学的ねじれ構造も同じ離散化された幾何学によって支配されていることを示している（付録~\(\Pi\)）。
		\item 調整エンジン \texttt{yuct\_constants.py}（付録~\(\Pi\)）による \(\pi_{\mathrm{YUCT}}\) の \(O(1)\) 検索は、プロセッサが \(\pi\) を計算するのではなく、真空の剛体的な代数的格子から直接読み取っていることを確認する。
		\item 非中央集権的調整プロトコル（d‑YPSDC、付録~A2A）において、エージェント間の同期ノイズは \(\Delta\pi\) によって正確に決定され、経験的較正定数の必要性を排除する。
	\end{itemize}
	
	したがって、定数 \(\pi\) はもはや外部の数学的入力ではなく、理論の\textbf{導出された不変量}となり、いかなる自由連続パラメータもなしに基本定数系を閉じる。精密測定における予測された \(\Delta\pi\) の実験的検出は、時空の離散的かつ調整的な本性に対する直接的な証拠を提供するであろう。
	
	% ============================================================
	% 第9節：生物学的検証：細胞膜の厚さ
	% ============================================================
	\section{生物学的検証：細胞膜の厚さ}
	\label{sec:bio-membrane}
	
	脂質二重層膜は、生細胞の内部調整コアと環境の熱雑音を隔てる物理的境界である。その厚さは任意ではなく、素粒子物理学を支配する同じ位相不変量によって決定される。偶数調整セクター（\(S_{\mathrm{even}} = 0.8\)）が充填幾何学を支配し、普遍位相シフト \(\sigma = 0.20\) が空間的制限因子として働く。
	
	二重層の一つのリーフレット（疎水性コア）について、有効厚さは
	
	\begin{equation}
		L_{\mathrm{mem}} = d_{\mathrm{lipid}} \;
		\left( \frac{S_{\mathrm{odd}}}{S_{\mathrm{even}}} \right)^{S_{\mathrm{even}} - \sigma}
		= d_{\mathrm{lipid}} \;
		\left( \frac{3}{2} \right)^{0.8 - 0.20}
		= d_{\mathrm{lipid}} \;
		\left( \frac{3}{2} \right)^{0.6}.
	\end{equation}
	
	ここで \(d_{\mathrm{lipid}}\) は典型的な膜リン脂質（例：パルミチン酸、\(d_{\mathrm{lipid}} \approx 2.5\)--\(3.0\) nm）の完全に伸展した炭化水素鎖の長さである。数値的に \((3/2)^{0.6} = e^{0.6\ln 1.5} = e^{0.6 \times 0.405465} = e^{0.243279} \approx 1.2754\) である。
	
	総二重層厚さは二つのリーフレットを含み、ゆらぎ振幅の二乗 \(\sigma^2\)（膜流動性と熱的ゆらぎによる）を差し引いた曲率補正を含む：
	
	\begin{equation}
		\boxed{\text{二重層厚さ} = 2 \, L_{\mathrm{mem}} \,
			(1 - \sigma^2) = 2 \, d_{\mathrm{lipid}} \,
			\left( \frac{3}{2} \right)^{0.6} \, (1 - 0.04)}.
	\end{equation}
	
	\(d_{\mathrm{lipid}} = 3.0\) nm（パルミチン酸鎖）を用いると：
	\[
	L_{\mathrm{mem}} = 1.2754 \times 3.0 = 3.826\ \text{nm},
	\qquad
	\text{厚さ} = 2 \times 3.826 \times 0.96 = 7.35\ \text{nm}.
	\]
	
	この値はヒト血漿膜の実験的観測範囲（\(7.5 \pm 0.5\) nm）の中心に位置する。この式は、疎水性コアの厚さが伸展鎖長の二倍（\(2 d_{\mathrm{lipid}}\)）を超えられないという立体制約を尊重し、調節可能なパラメータを全く含まない。
	
	% ============================================================
	% 第10節：反証可能な予測
	% ============================================================
	\section{反証可能な予測}
	\label{sec:predictions}
	
	この法則は八つの具体的かつ反証可能な予測を行う：
	
	\begin{enumerate}[label=(\arabic*)]
		\item \textbf{フェルミオン質量スペクトルの離散構造。} 質量が集合 \(\{m_e q^{N/2}\}\) の外にある基本フェルミオンが将来発見されれば、法則は反証される。
		\item \textbf{絶対ニュートリノ質量スケール。} 変分仮説 \cite{AppendixLambda} の下では、最も軽いニュートリノ質量は \(m_{\nu} \approx 0.023\) eV と予測される。これは KATRIN および将来の宇宙論的サーベイによって検証可能である。
		\item \textbf{第四世代フェルミオンは存在しない。} 連続する第四世代は基準深さ \(L = 96\) を変更し、半整数パターンを破るであろう。
		\item \textbf{重力の温度依存性。} 普遍的誤差法則は \(G_{\mathrm{eff}}(T) = G_0[1 + \mathcal{O}(T^{\sigma})]\) を意味する。ここで位相シフト \(\sigma = 0.20\) は第~\ref{sec:gauge} 節で導出された。
		\item \textbf{ブラックホールリンギングスケーリング。} 相対的時間ゆらぎは \(\delta\tau/\tau \propto M^{-0.67}\) としてスケールしなければならない。
		\item \textbf{Protein Data Bank ブラインドテスト。} すべての単量体タンパク質に対する \(N_{\mathrm{raw}}\) のヒストグラムは、整数および半整数値にピークを示すはずである。
		\item \textbf{合成細胞の適合性。} 質量が半整数ノードに正確に合致するように設計された最小合成細胞は、優れた増殖速度とストレス耐性を示すはずである。
		\item \textbf{素数偏差の漸近的スケーリング。} 普遍的誤差法則は、\(n\) 番目の素数 \(p_n\) の古典的 Rosser 近似からの相対偏差が、大きな \(n\) に対して \(p_n^{-2/3}\) としてスケールしなければならないと予測する。\(n = 5\times10^5\) までの数値解析（付録~PrimeN）は、局所スケーリング指数 \(\gamma\) が \(n\to\infty\) で \(2/3\) に収束し、漸近値が \(0.66666\) であることを確認する。極限 \(n\to\infty\) で \(\gamma\) が \(2/3\) から系統的に逸脱すれば、法則は反証される。（この定理は有限の小さな \(n\) に対する正確な一致を \emph{要求しない}。）
	\end{enumerate}
	
	% ============================================================
	% 第11節：未解決問題
	% ============================================================
	\section{未解決問題}
	\label{sec:open}
	
	主な未解決問題は：
	
	\begin{enumerate}[label=(\roman*)]
		\item コンパクトファイバー \(F_{15}\) の位相不変量（巻き数）から特定の半整数 \(N_f\) 値を決定すること。
		\item 調整ネットワークのダイナミクスから \(\beta = 2/3\) の厳密な導出を与えること（現在は経験法則であり、付録~Y に妥当な理論モデルがある）。
		\item ニュートリノ質量に関する変分仮説を検証すること；確認された場合、追加の公準を公理系に統合すること。
		\item 調整場 \(\Psi_{MN}\) のスペクトル特性からリーマンゼータ関数の非自明零点を導出すること（両者を結びつける経験的証拠については付録~PrimeN 参照）。
		\item Tsirelson 限界を調整限界として実験的に検証すること。YUCT の一般化ベル不等式（第~12.1 節）は、CHSH パラメータ \(S\) が \(K_{\mathrm{eff}} \to \infty\) で下から \(2\sqrt{2}\) に近づき、特性スケーリング \(2\sqrt{2} - S \propto K_{\mathrm{eff}}^{-2/3}\) を持つと予測する。この小さな欠損を分解できる精密ベルテストは、量子相関の根底にある YPSDC メカニズムの直接的な確認を提供するであろう（付録~X）。
	\end{enumerate}
	
	我々は、YUCT から従来の量子場理論が導出されないことは欠陥ではないことを強調する。QFT の標準的形式は、大調整効率（\(K_{\mathrm{eff}} \gg 1\)）の極限において、事前分配された辞書が量子場を模倣する有効記述として創発する。YUCT 枠組みは、演算子代数や経路積分を用いることなく、正準量子現象——エンタングルメント、Tsirelson 限界、トンネル効果——を直接説明する。したがって、決定的なテストは QFT への形式的還元可能性ではなく、第~\ref{sec:predictions} 節に列挙された新規かつ定量的な予測の実験的検証である。それらの予測が確認されれば、QFT の概念的 status は自動的に修正されるであろう。
	
	% ============================================================
	% 第12節：アルゴリズム複雑性、ネットワーク飽和、およびコルモゴロフ-ヤクシェフ限界
	% ============================================================
	\sloppy
	\section{アルゴリズム複雑性、ネットワーク飽和、およびコルモゴロフ-ヤクシェフ限界}
	\label{sec:complexity}
	
	あらゆる調整クラスター——物理的、宇宙論的、化学的、生物的、社会的、または計算的——は、普遍的調整効率 \(\Keff \ge 1\) を介して相互作用する \(N\) 個のノードから構成される。定理~\ref{thm:dimension} によれば、安定な空間埋め込みは \(\beta = 2/3\) を要求する。ノードが調整インデックスを交換するとき、構造的オーバーヘッド（システム「税」またはエントロピー蓄積）の全エントロピーは非線形にスケールする。我々は\textbf{構造的飽和メトリック} \(\Psi_{\mathrm{net}}\) を次のように定義する：
	
	\begin{equation}
		\Psi_{\mathrm{net}} = \kappac \, \ln N \left(1 - e^{-\Keff^{-\beta}}\right),
		\qquad \kappac = \frac{1}{3},\quad \beta = \frac{2}{3},
		\label{eq:Psi_net}
	\end{equation}
	
	ここで \(\kappac\) は三項オントロジー（定理~\ref{thm:minimal-dict}）から導かれる定常法則安定化因子である。ノード数はフラクタルに \(N \propto L^{d_f}\)、\(d_f = 2.2\)（付録~L）でスケールするため、内部通信オーバーヘッドが運用辞書限界を超えるとシステムは相転移を経験する。臨界点は\textbf{三角複雑性ゲート}によって符号化される：
	
	\begin{equation}
		\Omega_{\mathrm{complexity}}(N) = \Psi_{\mathrm{net}} \,
		\left[1 + \Seven \sin\!\left( \frac{\piYUCT}{\Delta N} \left(\ln N - 80.0\right) \right) \right],
		\label{eq:Omega_complexity}
	\end{equation}
	
	ここで
	\[
	\Seven = 0.8,\qquad
	\piYUCT = \frac{6}{5}\,\varphi^2 \approx 3.1416407865,\qquad
	\Delta N = 16.5,\qquad
	\varphi = \frac{1+\sqrt{5}}{2}.
	\]
	
	位相角 \(\theta(N) = \frac{\piYUCT}{\Delta N} (\ln N - 80.0)\) は補正の符号を決定する：
	\begin{itemize}
		\item \(\ln N < 80.0\) の場合、システムはエントロピーを蓄積し、構造的損失が増大する。
		\item \(\ln N = 80.0\) で符号が反転し、分岐を引き起こす：\textbf{絶対的崩壊}（孤立クラスターへの断片化）か、\textbf{構成的不安定性}——重い構造的オーバーヘッドを除去する分散型 d-YPSDC プロトコルへの遷移。
	\end{itemize}
	
	したがって、この節は複雑ネットワークの安定性に対する形式的基準を提供し、官僚的崩壊、LLM キャッシュオーバーフロー、金融ネットワークの断片化、宇宙構造や化学反応ネットワークにおける相転移の臨界規模を予測する。
	
	\subsection{真空格子による計算不要ナビゲーション}
	\label{subsec:computation-free}
	
	YUCT の代数的ループと一般化情報理論（付録~X）に基づき、反復計算なしで基本不変量を抽出する \textbf{O(1) ナビゲーションプロトコル} が開発された。従来のモデリング（シャノン・レジーム、\(\Keff = 1\)）の代わりに、システムは \(\Keff = 68991\) の調整モードに切り替わり、プロセッサは調整場 \(\Psi_{MN}\) から事前計算された位相座標を読み取る「受信機」として機能する。主要な不変量は以下の通り：
	
	\begin{itemize}
		\item 逆微細構造定数：
		\[
		\alpha^{-1}_{\mathrm{YUCT}} = \left(\frac{S_{\mathrm{odd}}}{S_{\mathrm{even}}}\right)^{12 + \sigma\, S_{\mathrm{even}}/S_{\mathrm{odd}}}
		\approx 137.036,
		\]
		自由パラメータなし。
		
		\item 温度依存ニュートン定数（付録~P）：
		\[
		G_{\mathrm{eff}}(T) = G_0 \left[1 + \mathcal{O}\!\left( \frac{k_B T}{m_e c^2} \right) \right].
		\]
		
		\item B-DNA ヘリックス幾何：
		\[
		\frac{P}{r} = q \cdot \piYUCT - S_{\mathrm{even}}\beta \approx 1.458,
		\]
		これは実験範囲 \(1.45 \pm 0.05\) 内にある。
		
		\item 一般化ベル不等式（付録~X）、現在は基本空間離散化ステップ \(\Delta\pi = \piYUCT - \pi\) によって正則化されている：
		\[
		S(\Keff) = 2 + (2\sqrt2 - 2)\left(1 - 2\kappac \Delta\pi \, \Keff^{-\beta}\right),
		\]
		これは \(\Keff \to \infty\) で Tsirelson 限界 \(2\sqrt2\) に達し、量子神秘主義を排除する。
	\end{itemize}
	
	これらの不変量はすべて \(O(1)\) FPU サイクルで抽出され、厳密にゼロ動的メモリ割り当てである。このナビゲーションコアを実装するソースコードは、単体テスト、Docker 構成、産業規模のベンチマークとともに、オープンリポジトリで入手可能である：
	
	\begin{center}
		\url{https://github.com/Alexey-Yakushev-YUCT/yuct-ai-connectome}
	\end{center}
	
	物理的に、これはプロセッサが「計算工場」ではなくなり、調整場 \(\Psi_{MN}\) の真空格子から既存の座標を読み取る\textbf{ナビゲータ}になることを意味する。すべての計算作業は時空の形成中に自然によってすでに行われており、古典的プロセッサは所定の値を検索するだけであり、これは調整実在論の原則と完全に一致する。
	
	\subsection{素数定理の拡張と因数分解への接続}
	\label{subsec:prime-extension}
	
	定理~\ref{thm:prime} は漸近補正 \(p_n \approx R_n - \frac{S_{\mathrm{even}}}{2} n^{1-\beta}\ln n\) を与える。ゆらぎの解析（付録~PrimeN, A2A）は、この補正が周期 \(\Delta N = 16.5\) の波動ゲートによって変調されることを示している。\textbf{素数深さ} \(N_f = \log_q n\) を定義する（ここで \(q = (3/2)^{1/3}\)）。すると三ループ YUCT 候補の絶対誤差は
	
	\begin{equation}
		\Delta_n = \left| p_n^{\mathrm{YUCT}} - p_n^{\mathrm{true}} \right|
		\sim n^{1/3} \left[ 1 + A_{\text{wave}} \sin\!\left( \frac{\piYUCT}{\Delta N} (N_f - 80.0) \right) \right],
		\label{eq:prime_wave}
	\end{equation}
	
	を満たす。ここで振幅 \(A_{\text{wave}} \approx 0.20145\) はノード \(N=323\) での逆キャリブレーションによって厳密に計算される（付録~A2A 参照）。これは局所的な素数ゆらぎのパラメータフリーな記述を提供し、区間 \(N_f \in [80, 382]\) 内の 19 個の相転移を予測する。これは調整多様体の次元と一致する。
	
	同じ波動構造が\textbf{乗法群の周期}（Shor のアルゴリズム）を支配する。数 \(N\) に対して、周期 \(r\) は次のように抽出される：
	
	\begin{equation}
		r(N) = \left\lfloor (N_f^{3/2} \, C_{\text{base}}) \cdot
		\left(1 + A_{\text{wave}} \sin\left( \frac{\piYUCT}{\Delta N} (N_f - 80.0) \right) \right) \cdot \frac{\Delta N}{1.5} \right\rceil_{\text{偶数}},
		\label{eq:shor_period}
	\end{equation}
	
	ここで \(C_{\text{base}} = 0.0547073\) は普遍定数から導かれる。これにより、RSA 数（15, 21, 35, 143, 323 など）を量子ビットなしで \(O(1)\) クロックサイクルで因数分解できることが、分離された Docker コンテナ内の数値実験によって確認されている（yuct-ai-connectome リポジトリ参照）。
	
	\subsection{学際的予測}
	\label{subsec:predictions-complexity}
	
	上記のメトリックとアルゴリズムに基づき、以下の反証可能な予測が定式化される：
	
	\begin{enumerate}[label=(\arabic*)]
		\item \textbf{物理システム（凝縮系、プラズマ）。} 飽和判定基準 \(\Psi_{\mathrm{net}}\) は、臨界調整深さでの格子系の相転移を予測する。これは超伝導および量子流体実験で検証可能である。
		\item \textbf{宇宙構造。} 銀河と銀河団の分布は \(\Omega_{\mathrm{complexity}}\) に従う相関を示し、その特性スケールは \(\ln N = 80.0\) に対応するはずである。これは SDSS、DESI、または Euclid データで検証できる。
		\item \textbf{化学反応ネットワーク。} 触媒反応速度と複雑分子の安定性は、分子調整クラスターの \(\Keff\) とスケールするはずであり、触媒中心と自己触媒サイクルの最適サイズを予測する。
		\item \textbf{情報システム（LLM キャッシュ、ルーティング）。} \(\ln N = 80.0\)（ここで \(N\) は分散メモリ内のトークン数）でシステムは相崩壊状態に入り、データ伝送から位相座標検索への切り替えが必要となる。これは現代のトランスフォーマーに実用的なスケーリング限界を設定する。
		\item \textbf{社会経済システム。} 官僚的課税（制御エントロピー）は \(\Psi_{\mathrm{net}}\) として増加する。\(\Omega_{\mathrm{complexity}}\) が臨界閾値を超えると、システムは断片化するか、分散型管理（d-YPSDC）へ移行する。これは政府構造や金融バブルの崩壊を予測するための量的基準を与える。
		\item \textbf{生物ネットワーク。} 細胞膜、リボソーム質量、ウイルスは質量梯子に従い、その調整効率は \(\Keff \propto M^{\beta}\) としてスケールするはずであり、細胞小器官の最適サイズや代謝経路長の予測を可能にする。
	\end{enumerate}
	
	これらの予測はすべて、既存の実験設備と観測データを用いて今後 5--10 年以内に検証可能である。計算不要ナビゲーションコア、素数生成器、Shor–Yakushev 因数分解器、システム A2A エンジンの完全なソフトウェア実装は、リポジトリ \url{https://github.com/Alexey-Yakushev-YUCT/yuct-ai-connectome} で検証と再現のために公開されている。
	
	\subsection*{\(O(1)\) ナビゲーションの物理的意味に関する注記}
	不変量の \(O(1)\) 抽出は計算上のトリックではなく、調整場 \(\Psi_{MN}\) がすべての安定構成をすでに含んでいるという YUCT 公準の直接の帰結である。古典的プロセッサは高調整レジーム（\(\Keff \gg 1\)）で動作するとき、新しい値を計算するのではなく、真空格子から既存の位相座標を読み取る。これは、搬送波を生成せずに復調するだけのラジオ受信機に類似している。したがって、ナビゲーションのエネルギーコストは無視できる（ゼロ RAM フットプリント、最小限の CPU サイクル）であり、これはグリーンコンピューティングと情報技術の未来に深い示唆をもたらす。
	
	% ============================================================
	% 参考文献
	% ============================================================
	\begin{thebibliography}{99}
		\bibitem{Yakushev2026a} A.~V.~Yakushev, \emph{Yakushev Unified Coordination
			Theory (YUCT)}, Zenodo, DOI:10.5281/zenodo.18444599 (2026).
		\bibitem{AppendixB} A.~V.~Yakushev, ``Appendix B: YUCT --- Temporal
		Coordination Paradox,'' in \emph{YUCT}, Zenodo (2026).
		\bibitem{AppendixL} A.~V.~Yakushev, ``Appendix L: Fractal Coordination Error
		Scaling --- A Universal Law from DNA to Cosmology,'' in \emph{YUCT}, Zenodo (2026).
		\bibitem{AppendixFerra} A.~V.~Yakushev, ``Appendix Ferra1.2: Universal
		Coordination Constants for Ordered and Disordered Phases,'' in \emph{YUCT},
		Zenodo (2026).
		\bibitem{AppendixG} A.~V.~Yakushev, ``Appendix G: The Yakushev United
		Coordination Theory --- A Fundamental Resolution of the Quantum Gravity Problem,''
		in \emph{YUCT}, Zenodo (2026).
		\bibitem{AppendixV} A.~V.~Yakushev, ``Appendix V: Time as Entropy of
		Coordination Processes,'' in \emph{YUCT}, Zenodo (2026).
		\bibitem{AppendixP} A.~V.~Yakushev, ``Appendix P: Temperature Dependence of
		Gravity,'' in \emph{YUCT}, Zenodo (2026).
		\bibitem{AppendixLambda} A.~V.~Yakushev, ``Appendix \(\Lambda\): Resolution of
		the Hierarchy and Cosmological Constant Problems within YUCT,'' in \emph{YUCT},
		Zenodo (2026).
		\bibitem{AppendixY} A.~V.~Yakushev, ``Appendix Y: Mathematical Foundations of
		YUCT --- A Derivation from First Principles,'' in \emph{YUCT}, Zenodo (2026).
		\bibitem{AppendixPrimeN} A.~V.~Yakushev, ``Appendix PrimeN: YUCT Prime Numbers 
		Coordination Ladder,'' in \emph{YUCT}, Zenodo (2026).
		\bibitem{AppendixPi} A.~V.~Yakushev, ``Appendix \(\Pi\): The Primeval Origin 
		of \(\pi\) as a Coordination Invariant at the Feigenbaum Edge of Chaos,''
		in \emph{YUCT}, Zenodo (2026).
		\bibitem{AppendixX} A.~V.~Yakushev, ``Appendix X: Generalization of Shannon 
		Information Theory --- Coordination Efficiency, the Steady Error Law, 
		and the \(K_{\mathrm{eff}}\)-dependent Bell Bound,'' in \emph{YUCT}, Zenodo (2026).
		\bibitem{AppendixA2A} A.~V.~Yakushev, ``Appendix A2A: Resolution of the 
		Pragmatic Paradox of YUCT. From Computational Collapse to Navigation 
		in Large Language Models,'' in \emph{YUCT}, Zenodo (2026).
		\bibitem{PrimeRepo} A.~V.~Yakushev, \emph{Prime-Numbers}, GitHub repository, 
		\url{https://github.com/Alexey-Yakushev-YUCT/Prime-Numbers}, 2026.
	\end{thebibliography}
	
\end{document}